\section[Sheet return, areal--volumetric incidence and constitutive orientation]{Sheet return, areal--volumetric incidence\\and constitutive orientation}
\label{iii:sec:hojas}

The sheet structure enters before electromagnetic quantities
are evaluated. Its function is to distinguish closure of a projection,
restoration of an orientation and evolution of the transported register.
These operations act on related objects of different types: a
projected trajectory may have closed while its lift remains
on the opposite sheet; orientation may have been restored while memory
preserves the two windings performed. The constitutive construction must receive
this information together with the structure determining it.

\subsection{Projected return and restoration of the oriented cover}

Let $\ell$ be a return loop of a surviving extension, and let
$\varepsilon:\langle\ell\rangle\to\{+1,-1\}$ be the character of its
lift, with $\varepsilon(\ell)=-1$. The associated interval bundle
has the presentation
\begin{equation}
 \mathcal M=([0,1]\times[-1,1])/{(0,u)\sim(1,-u)}.
 \label{iii:eq:mobius}
\end{equation}
The decisive datum is the sign of transport around the specified loop. The
identification of endpoints geometrically realizes that character.

\begin{proposition}[Two closure orders]
One winding of the loop closes the base and reverses the fibre orientation. Two
windings restore orientation. In the angular parametrization of the boundary
readout, these returns correspond, respectively, to the areal closure $2\pi$ and
the volumetric closure $4\pi$ of the oriented cover.
\end{proposition}
\begin{proof}
The equivalence relation of \eqref{iii:eq:mobius} identifies the terminal fibre
with the initial one by $u\mapsto-u$. Composing two transports gives
$u\mapsto-(-u)=u$. Parametrizing one winding of the base by $2\pi$ assigns $4\pi$
to its square. The statement concerns this lift and this readout.
\end{proof}

TPK memory also preserves the trajectory traversed and the register
updates. The square of the sign character is the identity, but it does not impose
vanishing of these updates. The spinorial representation, when
used, must transport this datum through its own map.

\subsection{Descent of incidence from the tritic calendar}

In the primitive block $(+1,-1,0)$, the TPK updates the arithmetic mode through
$+1:m\mapsto\Sigma$, $-1:m\mapsto\Pi$ and $0:m\mapsto m$. Its cyclic orbit
is $(\Sigma,\Pi,\Pi)$. The sheet readout assigns $\Sigma$ to the areal sheet and
$\Pi$ to the volumetric sheet. In the basis order $(e_{\rm ar},e_{\rm vol})$,
edge incidence is represented by
\begin{equation}
 F_{\rm av}=\begin{pmatrix}0&1\\1&1\end{pmatrix}.
 \label{iii:eq:fav}
\end{equation}
The three consecutive pairs are $\Sigma\to\Pi$, $\Pi\to\Pi$ and
$\Pi\to\Sigma$, each with multiplicity one. The calendars of lengths
$9$, $27$ and $108$ contain $3$, $9$ and $36$ copies of the primitive block.
Their incidence matrices are therefore these multiples of $F_{\rm av}$.

This construction preserves the distinction between chronological frequency and the measure
of the continuation space. The former assigns weights $(1/3,2/3)$ to the mode
cycle. The latter is defined on the minimal memory-one symbolic envelope
admitting exactly the preceding pairs. Its adjacency matrix is
$F_{\rm av}$; trajectory counts belong to that envelope.

\Needspace{10\baselineskip}
\begin{proposition}[Return weights of the two sheets]
Let $F_0=0$, $F_1=1$ and $F_{n+1}=F_n+F_{n-1}$. For $n\geq1$,
\begin{equation}
 F_{\rm av}^{n}=
 \begin{pmatrix}F_{n-1}&F_n\\F_n&F_{n+1}\end{pmatrix}.
 \label{iii:eq:fav-powers}
\end{equation}
The Parry measure of this envelope assigns the sheets weights proportional
to $(1,\varphi^2)$, and the ratio of diagonal counts converges to $\varphi^{-2}$.
\end{proposition}
\begin{proof}
The power formula follows by induction using the recurrence.
The square of $F_{\rm av}$ is strictly positive. Its positive Perron
vector is $(1,\varphi)$, with eigenvalue $\varphi$, since
$\varphi^2=\varphi+1$. As the matrix is symmetric, the Parry weights are
proportional to the squares of that vector's components. The other root
of the characteristic polynomial is $-\varphi^{-1}$; its modulus is strictly
less than $\varphi$, giving the limit of the diagonal ratios.
\end{proof}

Spectral identification of this incidence follows its construction
from the calendar. The regional coordinate $\varphi$ retains its earlier coinductive
genealogy: recognition of the eigenvalue does not select the calendar.

\subsection{Boundary marking and areal--volumetric ratio}

Let $\Pi_{\rm HMT}(X_\infty)$ be the value of the regional closure readout already
constructed. With the diagonal sheet projectors, define
\begin{equation}
 D_\pi=\Pi_{\rm HMT}(X_\infty)P_{\rm ar}+P_{\rm vol},\qquad
 \Theta_n=
 \frac{\langle e_{\rm ar},D_\pi F_{\rm av}^{n}e_{\rm ar}\rangle}
 {\langle e_{\rm vol},F_{\rm av}^{n}e_{\rm vol}\rangle}.
 \label{iii:eq:marked-return}
\end{equation}
Volumetric normalization is unity; regional closure marks the
areal return only once. Applying \eqref{iii:eq:fav-powers},
\begin{equation}
 \Theta_n=\Pi_{\rm HMT}(X_\infty)\frac{F_{n-1}}{F_{n+1}}
 \longrightarrow\frac{\pi}{\varphi^2}.
 \label{iii:eq:av-limit}
\end{equation}
The return ratio is determined by incidence, measure and marking.
The quotient $2\pi/4\pi$, by contrast, describes the cover's two closure orders.
Both relations participate in the sheet readout without replacing each other.

\subsection{Transport of orientation to the constitutive channels}

In a two-channel representation, let $\mathcal R$ be a self-adjoint
involution and $\mathcal S$ a unitary interchange operator satisfying
$\mathcal S\mathcal R\mathcal S^{-1}=-\mathcal R$. Write the angular coordinates
already constructed as $x=A_{\rm rad}$ and $y=C^*_{\rm rad}$, with $x>|y|$.
Then
\begin{equation}
 B=xI+y\mathcal R,\qquad T=e^{-B},\qquad
 \mathcal ST(y)\mathcal S^{-1}=T(-y).
 \label{iii:eq:channel-conjugation}
\end{equation}
Indeed, conjugation transforms $B(y)$ into $B(-y)$ and commutes with its convergent
exponential series. The eigenvalues $q_\pm=e^{-x\mp y}$ are interchanged.
To fix the response order, the chamber $x>y>0$ is used, where
$0<q_+<q_-<1$. The map $(x,y)\mapsto(x,-y)$ transports it to the
opposite chamber $x>-y>0$, where $0<q_-<q_+<1$. Both belong to the
contractive domain $x>|y|$. On the boundary $y=0$ the two channels coincide.
Conjugation interchanges the chambers; it does not preserve the strict ordering of
their labels. It acts on the projectors through
$\mathcal SP_\pm\mathcal S^{-1}=P_\mp$, with
$P_\pm=(I\pm\mathcal R)/2$.
This concrete map transmits orientation to the constitutive operator
developed next.
